%==============================================================================
% CAPÍTULO 3 — PYTHON PARA ODONTOLOGIA COMPUTACIONAL
% Google Colab, NumPy, Pandas e Visualização Científica
% Autor: Edson Laranjeiras
%==============================================================================

\chapter{Python para Odontologia Computacional}
\label{cap:python-colab}

\nivelzero

Todo profissional que deseja trabalhar com inteligência artificial odontológica precisa dominar as ferramentas computacionais que tornam a análise de dados clínicos e radiográficos possível. Este capítulo apresenta o ecossistema Python — NumPy, Pandas, Matplotlib e Seaborn — executado no ambiente Google Colab, que oferece acesso gratuito a GPU para treinamento de modelos. Mais do que um tutorial de programação, o capítulo estabelece a caixa de ferramentas essencial que será utilizada em todos os laboratórios práticos do livro, seguindo o protocolo SDD+TDD que garante reprodutibilidade e rigor científico a cada análise.

\subsection{Objetivos de Aprendizagem}

Ao final deste capítulo, você será capaz de:

\subsubsection{Configurar o Google Colab com GPU para Dados Odontológicos}
Aprenderá a ativar GPU Tesla T4 gratuita no Colab, montar o Google Drive e instalar as bibliotecas necessárias para processamento de imagens radiográficas.

\subsubsection{Manipular Dados Clínicos com NumPy e Pandas}
Representará radiografias como arrays NumPy e prontuários como DataFrames Pandas, realizando operações de filtragem, agregação e análise.

\subsubsection{Criar Visualizações Científicas com Matplotlib e Seaborn}
Construirá gráficos publicáveis -- histogramas, boxplots, curvas ROC -- para comunicar achados odontológicos de forma clara e profissional.

\subsubsection{Documentar Notebooks com SDD e TDD}
Aplicará o protocolo de desenvolvimento orientado a especificação e testes, garantindo que suas análises sejam verificáveis e reprodutíveis.

\subsubsection{Implementar Pipeline de Análise Exploratória com Testes}
Construirá um notebook completo que carrega, valida e visualiza um dataset real de radiografias, com testes automatizados em cada etapa.

\section{Fundamentação Teórica}

\subsection{Por que Python na Odontologia?}

Python \indice{Python} consolidou-se como a linguagem franca da ciência de dados e inteligência artificial por razões que transcendem modismos tecnológicos:

\begin{enumerate}
    \item \textbf{Ecossistema científico maduro.} Bibliotecas como NumPy (computação numérica), Pandas (análise de dados), Matplotlib/Seaborn (visualização), Scikit-learn (machine learning clássico), TensorFlow/PyTorch (deep learning) e OpenCV (processamento de imagens) formam um pipeline completo — da leitura de uma radiografia DICOM ao treinamento de uma CNN para diagnóstico.
    \item \textbf{Curva de aprendizado suave.} A sintaxe limpa e legível de Python permite que profissionais de saúde sem formação em ciência da computação se tornem produtivos em semanas, não meses.
    \item \textbf{Comunidade e reprodutibilidade.} Milhares de tutoriais, datasets públicos e modelos pré-treinados estão disponíveis. Notebooks Jupyter/Colab permitem documentar e compartilhar análises completas, promovendo a reprodutibilidade científica — um valor central deste livro.
    \item \textbf{Google Colab e acesso gratuito a GPU.} A barreira financeira para deep learning foi eliminada: GPUs NVIDIA Tesla T4, com 16 GB de VRAM, estão disponíveis gratuitamente no Colab, permitindo treinar modelos que há 5 anos exigiriam hardware de dezenas de milhares de reais.
\end{enumerate}

\subsection{Google Colab: Seu Laboratório de IA na Nuvem}

O Google Colaboratory (Colab) \indice{Google Colab} é um ambiente de notebooks Jupyter hospedado na nuvem que oferece:

\begin{itemize}
    \item \textbf{Execução sem instalação.} Navegador web + conta Google = ambiente de desenvolvimento completo.
    \item \textbf{GPU/TPU gratuita.} Menu Ambiente de Execução $\rightarrow$ Alterar tipo de ambiente de execução $\rightarrow$ GPU T4. Ideal para treinar CNNs em datasets de radiografias odontológicas.
    \item \textbf{Persistência no Google Drive.} Montagem do Drive permite salvar datasets, modelos treinados e resultados.
    \item \textbf{Células de código e texto.} Alternância entre explicações em Markdown/LaTeX e código executável em Python — o formato ideal para pesquisa reprodutível.
    \item \textbf{Integração com GitHub.} Salvar, versionar e compartilhar notebooks diretamente do Colab.
\end{itemize}

\subsection{Setup Inicial no Google Colab}

\begin{pratica}{Configuração do Google Colab para Odontologia Computacional}
\spec{
\textbf{Objetivo:} Configurar um notebook Colab com GPU e todas as bibliotecas necessárias para os laboratórios deste livro.

\textbf{Requisitos funcionais:}
\begin{itemize}
    \item Verificar disponibilidade e tipo de GPU
    \item Instalar e importar o pacote odontoia
    \item Montar Google Drive para persistência de dados
    \item Verificar versões de todas as bibliotecas
\end{itemize}

\textbf{Invariantes:}
\begin{itemize}
    \item O notebook deve funcionar em qualquer conta Google com GPU habilitada
    \item Todas as bibliotecas devem ser importadas sem erros
\end{itemize}
}
\end{pratica}

O primeiro passo é verificar o ambiente de execução:

\codigo{codigos/cap03/setup_colab_odontoia.py}

\teste{Verificação de ambiente: GPU Tesla T4 detectada com 15.0 GB VRAM. Pacote odontoia v0.1.0 instalado. Todas as dependências resolvidas. Google Drive montado em /content/drive.}

\section{Variáveis, Listas, Dicionários e Funções}

\subsection{Estruturas de Dados para o Cotidiano Odontológico}

Python oferece estruturas de dados nativas que mapeiam diretamente para conceitos odontológicos:

\begin{description}
    \item[Listas.] Sequências ordenadas e mutáveis. Ideal para armazenar números de dentes segundo notação FDI: \texttt{dentes\_posteriores = [16, 17, 26, 27, 36, 37, 46, 47]}.
    \item[Dicionários.] Pares chave-valor. Perfeito para mapear dente $\rightarrow$ condição clínica: \texttt{\{'16': 'cárie oclusal', '36': 'restauração insatisfatória'\}}.
    \item[Tuplas.] Sequências imutáveis. Útil para coordenadas de landmarks cefalométricos: \texttt{ponto\_A = (x, y)}.
    \item[Conjuntos (sets).] Coleções não ordenadas de elementos únicos. Para verificar quais superfícies dentárias foram restauradas: \texttt{superficies = \{'O', 'M', 'D', 'V', 'P'\}}.
\end{description}

\subsection{Funções para Lógica Clínica}

Funções permitem encapsular lógica clínica reutilizável. Por exemplo, a classificação de risco de cárie baseada no protocolo CAMBRA:

\begin{verbatim}
def classificar_risco_carie(cpob: int, fluxo_salivar: float,
                             frequencia_acucar: int, higiene: str) -> str:
    """Classifica risco de cárie segundo protocolo CAMBRA."""
    if cpob >= 6 or fluxo_salivar < 0.7:
        return "Alto Risco"
    elif cpob >= 3 or frequencia_acucar >= 4:
        return "Risco Moderado"
    else:
        return "Baixo Risco"
\end{verbatim}

\section{NumPy para Odontologia Computacional}

\subsection{O Array NumPy como Fundação}

\subsubsection{O que é o Array NumPy}

NumPy \indice{NumPy} (Numerical Python) é a biblioteca fundamental para computação científica em Python. Seu tipo central, o \texttt{ndarray} (N-dimensional array), é a estrutura sobre a qual praticamente todas as bibliotecas de IA são construídas.

\subsubsection{Aplicações em Imagens Odontológicas}

Para a odontologia computacional, arrays NumPy representam:
\begin{itemize}
    \item \textbf{Imagens radiográficas.} Uma radiografia panorâmica é um array 2D (altura $\times$ largura) de valores de intensidade (0 = preto, 255 = branco para 8 bits; 0--65535 para 16 bits DICOM).
    \item \textbf{Volumes CBCT.} Um exame CBCT é um array 3D (profundidade $\times$ altura $\times$ largura) onde cada voxel representa a densidade do tecido em unidades Hounsfield.
    \item \textbf{Matrizes de características.} Uma tabela com dados de 500 pacientes e 20 variáveis clínicas é um array 2D de forma \texttt{(500, 20)}.
    \item \textbf{Predições de modelos.} A saída de uma CNN para classificação de 10 tipos de lesões bucais em um lote de 32 imagens é um array de forma \texttt{(32, 10)}.
\end{itemize}

\subsection{Operações Essenciais com NumPy}

\begin{verbatim}
import numpy as np

# Criando arrays
radiografia = np.zeros((2048, 2048), dtype=np.uint16)  # Panorâmica 16-bit
cpod_scores = np.array([3, 8, 12, 1, 5, 9, 15, 2, 7, 11])

# Estatísticas descritivas de dados clínicos
media_cpod = np.mean(cpod_scores)        # 7.3
mediana_cpod = np.median(cpod_scores)    # 7.5
desvio_cpod = np.std(cpod_scores)        # 4.57

# Indexação e slicing para regiões de interesse (ROI)
# Extrair região periapical do dente 36 (coordenadas aproximadas)
roi_36 = radiografia[800:1000, 1200:1400]

# Operações vetorizadas (evitar loops em Python!)
# Normalizar valores de pixel para [0, 1]
radiografia_norm = (radiografia - radiografia.min()) / \
                   (radiografia.max() - radiografia.min())
\end{verbatim}

\subsection{Manipulando Imagens Odontológicas como Arrays}

Imagens radiográficas carregadas com OpenCV ou PIL são arrays NumPy. Isso significa que todas as operações de processamento de imagens — filtragem, equalização de histograma, detecção de bordas — podem ser implementadas como operações matriciais:

\begin{verbatim}
import cv2
import numpy as np

# Carregar radiografia panorâmica
pan = cv2.imread('panoramica_001.png', cv2.IMREAD_GRAYSCALE)
print(f"Forma: {pan.shape}, Tipo: {pan.dtype}")
# Output: Forma: (2408, 2992), Tipo: uint8

# Realce de contraste com equalização de histograma adaptativa (CLAHE)
# Essencial para visualizar detalhes sutis em radiografias odontológicas
clahe = cv2.createCLAHE(clipLimit=2.0, tileGridSize=(8,8))
pan_realcada = clahe.apply(pan)

# Detecção de bordas para identificar contornos dentários
bordas = cv2.Canny(pan_realcada, threshold1=50, threshold2=150)
\end{verbatim}

\section{Pandas para Dados Clínicos Odontológicos}

\subsection{O DataFrame como Prontuário Estruturado}

\subsubsection{O que é um DataFrame}

Pandas \indice{Pandas} é a biblioteca padrão para análise de dados tabulares. Seu tipo central, o \texttt{DataFrame}, é uma tabela bidimensional com colunas nomeadas — análogo a uma planilha Excel ou a uma tabela SQL.

\subsubsection{Aplicações Clínicas}

Na odontologia, DataFrames representam naturalmente:
\begin{itemize}
    \item Dados de prontuário eletrônico (paciente, idade, sexo, dentes tratados, diagnósticos)
    \item Resultados de exames periodontais (dente, face, profundidade de sondagem, sangramento, perda de inserção)
    \item Datasets de treinamento de machine learning (features clínicas + rótulos)
    \item Resultados de estudos clínicos (grupo, desfecho, covariáveis)
\end{itemize}

\subsection{Operações Essenciais com Pandas}

\begin{verbatim}
import pandas as pd

# Criar DataFrame com dados periodontais
df = pd.DataFrame({
    'dente': [16, 16, 16, 16, 16, 16, 11, 11, 11, 11],
    'face': ['MV', 'V', 'DV', 'ML', 'L', 'DL', 'MV', 'V', 'DV', 'ML'],
    'ps': [3, 2, 3, 5, 4, 3, 2, 2, 2, 3],      # Profundidade sondagem (mm)
    'ss': [1, 0, 1, 1, 1, 1, 0, 0, 0, 1],      # Sangramento à sondagem
    'nic': [0, 0, 0, 2, 1, 0, 0, 0, 0, 0],    # Nível de inserção clínica
    'placa': [1, 1, 1, 1, 0, 1, 0, 0, 0, 1],   # Índice de placa
    'paciente': ['P001'] * 10
})

# Análise descritiva
df.groupby('dente')['ps'].mean()     # PS média por dente
df[df['ps'] >= 4]                     # Filtrar sítios com bolsa >= 4mm
df['ss'].value_counts(normalize=True) # Proporção de sítios com sangramento

# Pivot table: PS média por dente e face
pivot = df.pivot_table(values='ps', index='dente',
                       columns='face', aggfunc='mean')
\end{verbatim}

\subsection{Importando Dados Odontológicos Reais}

\begin{verbatim}
from odontoia.data import load_panoramic

# Carregar dados do dataset público de radiografias panorâmicas
dataset = load_panoramic()
print(f"Imagens carregadas: {len(dataset)}")
print(f"Dimensões: {dataset[0].shape}")

# Criar DataFrame com metadados clínicos
metadados = pd.DataFrame({
    'id': range(len(dataset)),
    'tipo': ['panoramica'] * len(dataset),
    'dimensao_x': [img.shape[1] for img in dataset],
    'dimensao_y': [img.shape[0] for img in dataset]
})
metadados.describe()
\end{verbatim}

\section{Matplotlib e Seaborn para Visualização Científica}

\subsection{Fundamentos de Visualização Odontológica}

\subsubsection{A Importância da Visualização de Dados}

A visualização de dados não é mero adorno estético — é uma ferramenta analítica fundamental. Padrões invisíveis em tabelas de números tornam-se óbvios quando plotados. No contexto odontológico:

\begin{itemize}
    \item Um histograma de CPOD por faixa etária revela a distribuição de cárie na população;
    \item Um boxplot de profundidade de sondagem por sextante identifica padrões de doença periodontal;
    \item Uma curva ROC compara a performance de diferentes modelos de IA para diagnóstico de lesões bucais;
    \item Uma matriz de confusão visualiza os erros sistemáticos de um classificador de cáries.
\end{itemize}

\subsection{Matplotlib: Controle Absoluto}

Matplotlib \indice{Matplotlib} é a biblioteca de visualização fundamental do ecossistema Python. Oferece controle granular sobre cada elemento do gráfico:

\begin{verbatim}
import matplotlib.pyplot as plt
import numpy as np

# Distribuição de CPOD em uma população de 500 pacientes
np.random.seed(42)
idades = np.random.normal(35, 15, 500).clip(6, 90)
cpod = 0.5 * idades + np.random.normal(0, 4, 500)
cpod = cpod.clip(0, 32)

fig, ax = plt.subplots(1, 2, figsize=(12, 5))

# Histograma
ax[0].hist(cpod, bins=20, color='steelblue', edgecolor='white', alpha=0.8)
ax[0].set_xlabel('CPOD')
ax[0].set_ylabel('Frequência')
ax[0].set_title('Distribuição de CPOD na Amostra')
ax[0].axvline(np.mean(cpod), color='red', linestyle='--',
              label=f'Média = {np.mean(cpod):.1f}')
ax[0].legend()

# Scatter: CPOD vs Idade
ax[1].scatter(idades, cpod, alpha=0.5, c='steelblue', edgecolors='white')
ax[1].set_xlabel('Idade (anos)')
ax[1].set_ylabel('CPOD')
ax[1].set_title('CPOD vs Idade')

# Linha de tendência
z = np.polyfit(idades, cpod, 1)
p = np.poly1d(z)
x_trend = np.linspace(idades.min(), idades.max(), 100)
ax[1].plot(x_trend, p(x_trend), 'r--', linewidth=2,
           label=f'CPOD = {z[0]:.2f}*Idade + {z[1]:.1f}')
ax[1].legend()

plt.tight_layout()
plt.savefig('analise_cpod.png', dpi=300, bbox_inches='tight')
\end{verbatim}

\subsection{Seaborn: Elegância Estatística}

Seaborn \indice{Seaborn} é construído sobre Matplotlib e oferece uma API de alto nível otimizada para visualização de dados estatísticos. Seus pontos fortes incluem:

\begin{itemize}
    \item Integração nativa com DataFrames do Pandas;
    \item Paletas de cores otimizadas para percepção visual;
    \item Funções especializadas para distribuições, regressões, matrizes de correlação e dados categóricos;
    \item Estética ``pronta para publicação'' com poucas linhas de código.
\end{itemize}

\begin{verbatim}
import seaborn as sns
import pandas as pd

# Dados simulados de profundidade de sondagem por sextante
df_ps = pd.DataFrame({
    'Sextante': ['S1']*50 + ['S2']*50 + ['S3']*50 +
                ['S4']*50 + ['S5']*50 + ['S6']*50,
    'PS_mm': np.concatenate([
        np.random.normal(2.5, 0.8, 50),   # S1 - anterossuperior
        np.random.normal(2.3, 0.7, 50),   # S2 - póst-sup direito
        np.random.normal(2.4, 0.9, 50),   # S3 - póst-sup esquerdo
        np.random.normal(3.5, 1.2, 50),   # S4 - anteroinferior
        np.random.normal(3.8, 1.5, 50),   # S5 - póst-inf direito
        np.random.normal(4.0, 1.6, 50)    # S6 - póst-inf esquerdo
    ]).clip(1, 12),
    'Sangramento': np.random.binomial(1, 0.35, 300)
})

fig, axes = plt.subplots(1, 2, figsize=(14, 5))

# Boxplot da PS por sextante
sns.boxplot(data=df_ps, x='Sextante', y='PS_mm',
            palette='Blues', ax=axes[0])
axes[0].axhline(y=4, color='red', linestyle='--',
                label='Limiar periodontite (4mm)')
axes[0].set_title('Profundidade de Sondagem por Sextante')
axes[0].legend()

# Violin plot com distribuição
sns.violinplot(data=df_ps, x='Sextante', y='PS_mm',
               palette='Blues', inner='quartile', ax=axes[1])
axes[1].set_title('Distribuição da PS por Sextante (Violin Plot)')
plt.tight_layout()
\end{verbatim}

\section{Documentação de Notebooks Científicos}

\subsection{O Padrão SDD+TDD para Notebooks Odontológicos}

Neste livro, todo notebook segue o protocolo SDD (Specification-Driven Development) + TDD (Test-Driven Development) \indice{SDD+TDD}:

\begin{enumerate}
    \item \textbf{Especificação (SDD).} Antes de escrever uma linha de código, definimos formalmente: objetivo, dados de entrada, saída esperada, invariantes que devem ser preservados e critérios de aceitação.
    \item \textbf{Testes (TDD).} Para cada função, escrevemos primeiro o teste que define o comportamento esperado. O código só é considerado completo quando todos os testes passam.
    \item \textbf{Documentação inline.} Cada célula de código é precedida por uma célula Markdown explicando: o que faz, por que é necessário, quais conceitos odontológicos são aplicados.
    \item \textbf{Reprodutibilidade.} O notebook inclui seeds fixas, versões de bibliotecas documentadas e caminhos de dados relativos.
\end{enumerate}

\subsection{Estrutura de um Notebook Científico Odontológico}

\begin{verbatim}
# ============================================================
# TÍTULO: Análise Exploratória de Dados Periodontais
# AUTOR: Edson Laranjeiras
# DATA: 2026-07-09
# OBJETIVO: Identificar padrões de doença periodontal em
#           uma coorte de 300 pacientes.
# INVARIANTES:
#  - Dados não devem conter valores negativos de PS
#  - Cada paciente deve ter no máximo 168 sítios (28 dentes x 6 faces)
#  - A soma de PS deve ser monotônica com a idade
# ESPECIFICAÇÃO SDD: SPEC-PERIO-001
# ============================================================

# Célula 1: Configuração e imports
import numpy as np
import pandas as pd
import matplotlib.pyplot as plt
import seaborn as sns
from odontoia.data import load_periodontal_cohort

# Célula 2: Carregamento e validação (TDD)
def test_data_integrity(df):
    """TESTE TDD: Verifica invariantes do dataset."""
    assert (df['ps_mm'] >= 0).all(), "PS negativa detectada"
    assert df.groupby('paciente_id').size().max() <= 168, \
           "Paciente com mais de 168 sítios"
    assert df['ps_mm'].max() <= 15, "PS > 15mm — outlier suspeito"
    print("TESTE PASSOU: Integridade dos dados verificada.")

df = load_periodontal_cohort()
test_data_integrity(df)
\end{verbatim}

\section{Revisão Bibliográfica Auditável}

O ecossistema Python para ciência de dados e IA odontológica é sustentado por bibliotecas com ampla documentação e comunidades ativas. A Tabela~\ref{tab:bibliotecas-python} apresenta as principais bibliotecas utilizadas neste livro.

\begin{table}[H]
  \centering
  \caption{Bibliotecas Python para Odontologia Computacional}
  \label{tab:bibliotecas-python}
  \begin{tabularx}{\textwidth}{lXcc}
    \toprule
    \textbf{Biblioteca} & \textbf{Função na Odontologia Computacional} & \textbf{Desde} & \textbf{Substitui} \\
    \midrule
    NumPy & Arrays para imagens radiográficas e operações matriciais & 2006 & MATLAB (parcial) \\
    Pandas & Manipulação de dados clínicos tabulares (prontuários, exames) & 2008 & Excel, SQL \\
    Matplotlib & Visualização científica de dados odontológicos & 2003 & MATLAB plotting \\
    Seaborn & Visualização estatística de alto nível & 2012 & — \\
    OpenCV & Processamento de imagens radiográficas (CLAHE, filtros, bordas) & 2000 & — \\
    Scikit-learn & ML clássico para predição de risco e classificação & 2007 & — \\
    TensorFlow & Deep learning para CNNs odontológicas (detecção, segmentação) & 2015 & — \\
    PyTorch & Deep learning com grafos dinâmicos (pesquisa odontológica) & 2016 & — \\
    MONAI & Framework de IA para imagem médica (DICOM, transformações 3D) & 2020 & — \\
    odontoia & Pacote específico para datasets, métricas e modelos odontológicos & 2026 & — \\
    \bottomrule
  \end{tabularx}
\end{table}

\section{Notas de Rodapé Explicativas}

\notaexplicativa{GPU (Graphics Processing Unit, ou Unidade de Processamento Gráfico) é um processador especializado originalmente projetado para renderização gráfica, mas que se mostrou excepcionalmente eficiente para as operações de álgebra linear que constituem o núcleo do deep learning. Enquanto uma CPU típica possui 4--16 núcleos otimizados para processamento sequencial, uma GPU moderna possui milhares de núcleos otimizados para processamento paralelo — ideal para multiplicar matrizes de pesos por milhões de pixels simultaneamente.}

\notaexplicativa{DICOM (Digital Imaging and Communications in Medicine) é o padrão internacional (ISO 12052) para armazenamento e transmissão de imagens médicas. Arquivos DICOM (.dcm) contêm não apenas os pixels/voxels da imagem, mas também metadados essenciais: identificação do paciente, parâmetros de aquisição (kVp, mAs, corrente do tubo), orientação espacial e informações do equipamento. Radiografias odontológicas digitais e exames CBCT são armazenados neste formato.}

\notaexplicativa{CLAHE (Contrast Limited Adaptive Histogram Equalization) é uma técnica de processamento de imagens que melhora o contraste local sem amplificar ruído. Em radiografias odontológicas, o CLAHE é particularmente útil para realçar detalhes sutis como lesões de cárie incipientes ou trabeculado ósseo alveolar fino, que podem ser pouco visíveis em imagens com contraste global inadequado.}

\section{Laboratório em Python e Google Colab}

\subsection{Notebook Completo: Análise Exploratória de Radiografias Panorâmicas}

\begin{pratica}
\spec{
\textbf{SPEC-COL-001: Análise Exploratória de Dataset de Radiografias Panorâmicas}

\textbf{Objetivo:} Construir um notebook completo de análise exploratória de dados (AED) para um dataset público de radiografias panorâmicas odontológicas, aplicando SDD e TDD.

\textbf{Requisitos funcionais:}
\begin{itemize}
    \item Carregar dataset usando \texttt{odontoia.data.load\_panoramic()}
    \item Calcular estatísticas descritivas: número de imagens, dimensões médias, distribuição de intensidade de pixels
    \item Visualizar distribuição de dimensões (histograma), intensidade média por imagem (boxplot)
    \item Exibir 9 exemplos do dataset em grid $3 \times 3$ com anotações
    \item Identificar e reportar outliers (imagens com dimensões ou intensidade anômalas)
    \item Salvar relatório automático em PDF
    \item Todos os testes TDD devem passar
\end{itemize}

\textbf{Invariantes:}
\begin{itemize}
    \item Todas as imagens devem ser arrays NumPy 2D com dtype uint8 ou uint16
    \item Nenhuma imagem deve ter dimensão zero em qualquer eixo
    \item A seed aleatória deve ser fixada para reprodutibilidade
\end{itemize}

\textbf{Critérios de aceitação:}
\begin{itemize}
    \item Testes TDD verificam carregamento, dimensões, tipo e ausência de NaNs
    \item Visualizações são salvas em PNG (300 dpi)
    \item Relatório inclui: estatísticas descritivas, gráficos e lista de outliers
    \item Código documentado com docstrings e comentários em português
\end{itemize}
}
\end{pratica}

\subsection{Código do Notebook}

\codigo{codigos/cap03/notebook_aed_panoramicas.py}

\subsection{Testes TDD}

\codigo{codigos/cap03/testes_aed_panoramicas.py}

\teste{Todos os 8 testes passaram: (1) carregamento do dataset — 350 imagens carregadas, (2) verificação de tipo — todas uint8 ou uint16, (3) verificação de dimensões — sem imagens com dimensão zero, (4) verificação de NaNs — zero valores nulos, (5) estatísticas descritivas — média de dimensão $2014 \times 2956$ pixels, (6) detecção de outliers — 2 imagens identificadas com desvio $>3\sigma$, (7) grid $3 \times 3$ salvo com sucesso, (8) relatório PDF gerado com 12 páginas.}

\subsection{Interpretação do Laboratório}

A análise exploratória revelou padrões importantes:
\begin{enumerate}
    \item \textbf{Variabilidade dimensional.} As radiografias panorâmicas variam de $1890 \times 1420$ a $3024 \times 4096$ pixels — uma variação de 2.3$\times$ em área. Modelos de deep learning devem incluir redimensionamento padronizado como etapa de pré-processamento.
    \item \textbf{Distribuição de intensidade.} Imagens apresentam modos bimodais no histograma: um pico para tecidos moles/ar (valores baixos) e outro para estruturas mineralizadas (valores altos). Este padrão é esperado e desejável para segmentação.
    \item \textbf{Outliers.} As 2 imagens identificadas como outliers apresentavam superexposição — saturação em mais de 40\% dos pixels, tornando-as inadequadas para treinamento diagnóstico. Sua exclusão do dataset é justificada.
\end{enumerate}

\section{Síntese do Capítulo}

Neste capítulo, estabelecemos a fundação computacional sobre a qual todos os laboratórios subsequentes do livro serão construídos. Python, através do ecossistema NumPy/Pandas/Matplotlib/Seaborn, oferece ao cirurgião-dentista e ao pesquisador odontológico as ferramentas para:

\begin{itemize}
    \item Representar imagens radiográficas e dados clínicos como estruturas de dados computacionalmente eficientes;
    \item Realizar análises estatísticas e visuais que revelam padrões invisíveis a olho nu;
    \item Documentar todo o processo analítico de forma reprodutível, usando o protocolo SDD+TDD que será aprofundado nos capítulos da Parte III.
\end{itemize}

O Google Colab democratiza o acesso à infraestrutura computacional necessária para deep learning, eliminando a barreira financeira que historicamente limitava a pesquisa em IA odontológica a instituições com recursos computacionais substanciais.

\section{OpenCV e Processamento de Imagens Odontológicas}

\subsection{Por que OpenCV na Odontologia?}

\subsubsection{O que é o OpenCV}

OpenCV (Open Source Computer Vision) \indice{OpenCV} é a biblioteca padrão para processamento de imagens e visão computacional. Em odontologia, é utilizada para:

\begin{itemize}
    \item \textbf{Pré-processamento de radiografias.} Redimensionamento, normalização, equalização de histograma e CLAHE preparam imagens para CNNs.
    \item \textbf{Detecção de landmarks.} Identificação automatizada de pontos anatômicos em cefalométricas.
    \item \textbf{Segmentação básica.} Limiarização (thresholding) e operações morfológicas para isolar estruturas dentárias.
    \item \textbf{Augmentação de dados.} Rotação, translação, espelhamento e ajuste de brilho/contraste para aumentar datasets de treinamento.
    \item \textbf{Análise de textura.} Extração de features como GLCM (Gray-Level Co-occurrence Matrix) para radiomics.
\end{itemize}

\subsection{Operações Essenciais com OpenCV para Radiografias}

\begin{verbatim}
import cv2
import numpy as np

# 1. Carregamento e conversão
# Radiografias odontológicas frequentemente vêm em 16-bit DICOM
pan = cv2.imread('panoramica.dcm', cv2.IMREAD_UNCHANGED)
if pan.dtype == np.uint16:
    # Normalizar 16-bit -> 8-bit para visualização
    pan_8bit = cv2.normalize(pan, None, 0, 255,
                             cv2.NORM_MINMAX, dtype=cv2.CV_8U)

# 2. CLAHE — realce de contraste local
clahe = cv2.createCLAHE(clipLimit=2.0, tileGridSize=(8,8))
pan_clahe = clahe.apply(pan_8bit)

# 3. Redução de ruído (preservando bordas)
pan_denoised = cv2.bilateralFilter(pan_clahe, d=9,
                                    sigmaColor=75, sigmaSpace=75)

# 4. Limiarização adaptativa para detectar dentes
thresh = cv2.adaptiveThreshold(
    pan_denoised, 255,
    cv2.ADAPTIVE_THRESH_GAUSSIAN_C,
    cv2.THRESH_BINARY_INV, 11, 2
)

# 5. Operações morfológicas para limpar a segmentação
kernel = np.ones((3,3), np.uint8)
morph = cv2.morphologyEx(thresh, cv2.MORPH_CLOSE, kernel, iterations=2)
morph = cv2.morphologyEx(morph, cv2.MORPH_OPEN, kernel, iterations=1)

# 6. Encontrar contornos (potenciais dentes)
contours, _ = cv2.findContours(morph, cv2.RETR_EXTERNAL,
                                cv2.CHAIN_APPROX_SIMPLE)

# Filtrar contornos por área (remover ruído)
min_area = 500  # pixels — ajustar conforme resolução
dentes = [c for c in contours if cv2.contourArea(c) > min_area]
print(f"Contornos detectados: {len(dentes)} possíveis dentes")
\end{verbatim}

\subsection{Manipulando Formato DICOM com pydicom}

Arquivos DICOM (.dcm) requerem bibliotecas especializadas:

\begin{verbatim}
import pydicom
import matplotlib.pyplot as plt

# Carregar arquivo DICOM de CBCT ou panorâmica
ds = pydicom.dcmread('exame_cbct.dcm')

# Extrair metadados clínicos
print(f"Paciente: {ds.PatientName}")
print(f"Data do exame: {ds.StudyDate}")
print(f"Modalidade: {ds.Modality}")
print(f"kVp: {ds.KVP}")        # Tensão do tubo
print(f"mAs: {ds.XRayTubeCurrent}")  # Corrente

# Acessar os pixels
pixels = ds.pixel_array  # Array NumPy!
print(f"Dimensões: {pixels.shape}")

# Para CBCT 3D, pixels.shape pode ser (N_slices, H, W)
if len(pixels.shape) == 3:
    slice_central = pixels[pixels.shape[0] // 2]
    plt.imshow(slice_central, cmap='gray')
    plt.title(f'Slice central do CBCT — {ds.PatientName}')
    plt.axis('off')
    plt.show()

# Window/Level para tecido ósseo (Unidades Hounsfield)
# Osso: WL=400, WW=1500
window_center = 400
window_width = 1500
vmin = window_center - window_width // 2
vmax = window_center + window_width // 2
plt.imshow(slice_central, cmap='gray', vmin=vmin, vmax=vmax)
\end{verbatim}

\section{Manipulação Avançada de Dados Clínicos com Pandas}

\subsection{Limpando e Preparando Dados Odontológicos}

Dados clínicos reais raramente chegam ``limpos''. O Pandas oferece ferramentas para lidar com valores ausentes, inconsistências e outliers:

\begin{verbatim}
import pandas as pd
import numpy as np

# Carregar dataset com problemas típicos
df = pd.read_csv('dados_periodontais.csv')

# 1. Diagnóstico de qualidade dos dados
print(f"Registros: {len(df)}")
print(f"Valores nulos:\n{df.isnull().sum()}")
print(f"Tipos de dados:\n{df.dtypes}")

# 2. Tratar valores ausentes
# PS (Profundidade de Sondagem) — inputar com mediana
df['ps_mm'] = df['ps_mm'].fillna(df['ps_mm'].median())

# 3. Detectar e tratar outliers
# PS > 15mm é fisiologicamente implausível
outliers_ps = df[df['ps_mm'] > 15]
print(f"Outliers de PS removidos: {len(outliers_ps)}")
df = df[df['ps_mm'] <= 15]

# 4. Criar features derivadas
# Classificar sítios por severidade
df['severidade'] = pd.cut(
    df['ps_mm'],
    bins=[0, 3, 5, 7, 20],
    labels=['Saúde', 'Leve', 'Moderada', 'Severa']
)

# Média de PS por dente para cada paciente
ps_por_dente = df.groupby(['paciente_id', 'dente'])['ps_mm'].agg(
    ['mean', 'max', 'std', 'count']
).reset_index()

# 5. Encoding para ML: one-hot para variáveis categóricas
df_encoded = pd.get_dummies(
    df,
    columns=['severidade', 'face', 'sangramento'],
    drop_first=True  # Evitar armadilha da variável dummy
)
\end{verbatim}

\subsection{Merge e Join de Dados Multimodais}

Na odontologia, dados de um paciente frequentemente estão fragmentados em múltiplas fontes:

\begin{verbatim}
# Dados demográficos
demog = pd.DataFrame({
    'paciente_id': ['P001', 'P002', 'P003'],
    'idade': [45, 32, 58],
    'sexo': ['F', 'M', 'F'],
    'diabetes': [False, False, True],
    'tabagismo': ['nunca', 'atual', 'ex-fumante']
})

# Dados periodontais (múltiplos registros por paciente)
perio = ps_por_dente  # DataFrame criado anteriormente

# Dados radiográficos (features extraídas por CNN)
radio_features = pd.DataFrame({
    'paciente_id': ['P001', 'P002', 'P003'],
    'perda_ossea_mm': [1.2, 3.4, 5.1],
    'furca_score': [0, 1, 3]
})

# Merge: consolidar múltiplas fontes
df_completo = demog.merge(radio_features, on='paciente_id')
df_completo['ps_media'] = perio.groupby('paciente_id')['mean'].mean()
df_completo['ps_max'] = perio.groupby('paciente_id')['max'].max()

# Dataset pronto para modelagem preditiva
print(df_completo.head())
\end{verbatim}

\section{Visualização Avançada para Publicação Científica}

\subsection{Estilo Publicável com Matplotlib}

\begin{verbatim}
import matplotlib.pyplot as plt
import seaborn as sns

# Configurar estilo para publicação
plt.style.use('seaborn-v0_8-whitegrid')
plt.rcParams.update({
    'font.family': 'serif',
    'font.size': 11,
    'axes.titlesize': 13,
    'axes.labelsize': 12,
    'figure.dpi': 150,
    'savefig.dpi': 300,
    'savefig.bbox': 'tight',
    'savefig.format': 'png'
})

# Criar figura multi-painel no estilo de periódico científico
fig = plt.figure(figsize=(14, 10))
gs = fig.add_gridspec(2, 3, hspace=0.35, wspace=0.35)

# Painel A: Distribuição de CPOD por faixa etária
ax1 = fig.add_subplot(gs[0, 0])
# ... (código de plotagem)

# Painel B: Boxplot de PS por sextante com significância estatística
ax2 = fig.add_subplot(gs[0, 1])
# ... (código de plotagem)

# Painel C: Curva ROC com IC95% (bootstrap)
ax3 = fig.add_subplot(gs[0, 2])
# ... (código de plotagem)

# Painel D: Matriz de correlação (heatmap)
ax4 = fig.add_subplot(gs[1, 0])
# ... (código de plotagem)

# Painel E: Scatter com linha de regressão
ax5 = fig.add_subplot(gs[1, 1])
# ... (código de plotagem)

# Painel F: Gráfico radar de performance de múltiplos modelos
ax6 = fig.add_subplot(gs[1, 2], projection='polar')
# ... (código de plotagem)

# Adicionar labels (A)-(F) para referência no texto
for ax, label in zip([ax1, ax2, ax3, ax4, ax5, ax6],
                      ['A', 'B', 'C', 'D', 'E', 'F']):
    ax.text(-0.1, 1.15, f'({label})', transform=ax.transAxes,
            fontsize=14, fontweight='bold')

plt.savefig('figura_publicacao.png', dpi=300, bbox_inches='tight')
\end{verbatim}

\section{Integração com o Pacote odontoia}

\subsection{Visão Geral do odontoia}

O pacote \texttt{odontoia} \indice{odontoia (pacote)} foi desenvolvido especificamente para este livro e consolida funções reutilizáveis para pesquisa em IA odontológica. Principais submódulos:

\begin{description}
    \item[\texttt{odontoia.data}.] Carregadores de datasets públicos (panorâmicas, CBCT, periodontais) com validação integrada.
    \item[\texttt{odontoia.metrics}.] Calculadoras de métricas diagnósticas (sensibilidade, especificidade, AUC, Dice) com intervalos de confiança.
    \item[\texttt{odontoia.models}.] Wrappers para modelos de ML clássico (\texttt{dental\_ml}) e deep learning (\texttt{dental\_dl}) com interfaces padronizadas.
    \item[\texttt{odontoia.viz}.] Visualizações especializadas: odontogramas interativos, curvas ROC comparativas, mapas de ativação Grad-CAM.
    \item[\texttt{odontoia.preprocessing}.] Funções de pré-processamento odontológico: CLAHE, window/level DICOM, extração de ROI dentária.
\end{description}

\subsection{Exemplos de Uso do odontoia}

\begin{verbatim}
from odontoia.data import load_panoramic, load_periodontal_cohort
from odontoia.metrics import classification_report_ci
from odontoia.viz import plot_roc_comparison, plot_odontogram
from odontoia.preprocessing import apply_clahe, extract_tooth_roi

# Carregar e pré-processar
images = load_panoramic()
processed = [apply_clahe(img) for img in images]

# Extrair ROI do dente 36 (coordenadas FDI)
roi_36 = [extract_tooth_roi(img, tooth=36) for img in processed]

# Métricas com intervalos de confiança
y_true = [...]  # Ground truth
y_pred = [...]  # Predições do modelo
report = classification_report_ci(y_true, y_pred)
print(report)

# Visualização especializada
plot_odontogram(paciente_id='P001', dados_periodontais=df_perio)
plot_roc_comparison(
    modelos={'XGBoost': y_pred_xgb, 'RandomForest': y_pred_rf},
    y_true=y_true
)
\end{verbatim}

\section{Boas Práticas para Notebooks Odontológicos}

\subsection{Estrutura Recomendada}

Todo notebook científico odontológico deve seguir esta estrutura:

\begin{enumerate}
    \item \textbf{Cabeçalho padronizado.} Título, autor, data, objetivo, especificação SDD (SPEC-ID), invariantes.
    \item \textbf{Setup.} Imports, configuração de seeds, verificação de ambiente (GPU disponível?).
    \item \textbf{Carregamento de dados.} Com validação automática (testes TDD inline).
    \item \textbf{Análise exploratória.} Estatísticas descritivas, visualizações, detecção de outliers.
    \item \textbf{Pré-processamento.} Limpeza, normalização, encoding, particionamento.
    \item \textbf{Modelagem.} Treinamento com validação cruzada, otimização de hiperparâmetros.
    \item \textbf{Avaliação.} Métricas com IC95\%, curvas ROC, matriz de confusão, análise de erros.
    \item \textbf{Interpretação.} Feature importance, SHAP, análise de casos discordantes.
    \item \textbf{Conclusão.} Síntese dos achados, limitações, próximos passos.
    \item \textbf{Exportação.} Salvar modelo, figuras e relatório.
\end{enumerate}

\subsection{Checklist de Qualidade para Notebooks}

\begin{itemize}
    \item \textbf{Reprodutibilidade:} Seeds fixas (\texttt{np.random.seed(42)}), versões de bibliotecas documentadas
    \item \textbf{Testabilidade:} Funções com docstrings, testes TDD para invariantes e casos-limite
    \item \textbf{Legibilidade:} Células Markdown explicativas entre células de código; nomes de variáveis em português claro
    \item \textbf{Eficiência:} Evitar loops Python em dados grandes; preferir operações vetorizadas do NumPy
    \item \textbf{Portabilidade:} Caminhos relativos; dependências declaradas em \texttt{requirements.txt}
    \item \textbf{Segurança:} NUNCA hardcodear credenciais, tokens ou caminhos absolutos locais
\end{itemize}

\section{Colab Pro, Colab Pro+ e Alternativas}

\subsection{Além da GPU Gratuita}

Embora a GPU gratuita do Colab (Tesla T4 com 15 GB VRAM) seja suficiente para a maioria dos experimentos deste livro, projetos mais ambiciosos podem exigir mais recursos:

\begin{description}
    \item[Colab Pro (US\$ 9.99/mês).] GPUs mais rápidas (Tesla V100, A100), mais memória RAM (25 GB vs. 12 GB), execuções mais longas (24h vs. 12h) e menos desconexões por inatividade.
    \item[Colab Pro+ (US\$ 49.99/mês).] Execução em segundo plano (fechar o navegador não interrompe), ainda mais RAM (52 GB), acesso prioritário a GPUs.
    \item[Kaggle Notebooks (gratuito).] Alternativa ao Colab com 30h/semana de GPU Tesla P100. Interface similar Jupyter.
    \item[Google Cloud Vertex AI.]. Para projetos profissionais que exigem pipelines de treinamento escaláveis, versionamento de modelos e deploy em produção.
\end{description}

Para a grande maioria dos leitores, o Colab gratuito é mais que suficiente para todos os 93 laboratórios e projetos guiados deste livro.

\subsection{Alternativas Locais}

Se você prefere desenvolver localmente:

\begin{description}
    \item[Anaconda Distribution.] Gerenciador de ambientes Python com instalação one-click de NumPy, Pandas, Jupyter e 250+ pacotes científicos. Gratuito.
    \item[VS Code + Jupyter extension.] Editor de código profissional com suporte nativo a notebooks. Gratuito.
    \item[PyCharm Professional.] IDE completa com suporte a notebooks, debugging remoto e integração com bancos de dados. Versão educacional gratuita para estudantes e professores.
\end{description}

\section{Depurando Código Python para Odontologia}

\subsection{Erros Comuns e Como Evitá-los}

\begin{description}
    \item[IndexError em DICOM 3D.] CBCTs são arrays 3D (slice, altura, largura). Acessar \texttt{pixels[y, x]} quando deveria ser \texttt{pixels[slice\_idx, y, x]} é o erro mais comum.
    \item[MemoryError com datasets grandes.] Carregar 500 panorâmicas de 3000$\times$4000 pixels em float32 consome $\approx 500 \times 3000 \times 4000 \times 4$ bytes $\approx 24$ GB — mais que a RAM disponível. Solução: processamento em batches ou redução de resolução.
    \item[Shape mismatch em CNNs.] Modelos de classificação esperam entrada (batch, canais, altura, largura). Se sua imagem é (altura, largura), adicione dimensões: \texttt{np.expand\_dims(img, axis=(0, 1))}.
    \item[Dtype mismatch.] Radiografias 16-bit (0--65535) alimentadas em modelos treinados com 8-bit (0--255) produzem predições incorretas. Sempre verifique e normalize: \texttt{img = img.astype(np.float32) / img.max()}.
    \item[ValueError em Pandas merge.] Colunas com mesmo nome mas tipos diferentes (ex.: \texttt{paciente\_id} como int em um DataFrame e string em outro) causam falhas silenciosas no merge.
\end{description}

\subsection{Ferramentas de Depuração}

\begin{itemize}
    \item \textbf{assert.} Verificações inline para invariantes: \texttt{assert img.shape == (2408, 2992), f"Forma inesperada: \{img.shape\}"}
    \item \textbf{pdb / ipdb.} Debugger interativo para inspecionar variáveis e executar passo a passo.
    \item \textbf{\%debug magic (IPython).} Após uma exceção, digite \texttt{\%debug} no notebook para entrar no debugger no ponto exato do erro — recurso subestimado e extremamente útil.
    \item \textbf{tqdm.} Barras de progresso para loops longos, essencial para processamento de datasets: \texttt{for img in tqdm(images, desc="Processando"): ...}
\end{itemize}

\section{Leituras Recomendadas}
\begin{enumerate}
    \item VanderPlas J. (2016). \textit{Python Data Science Handbook}. O'Reilly Media. — Referência essencial para NumPy, Pandas e Matplotlib.
    \item McKinney W. (2017). \textit{Python for Data Analysis} (2nd ed.). O'Reilly Media. — Escrito pelo criador do Pandas.
    \item Géron A. (2022). \textit{Hands-On Machine Learning with Scikit-Learn, Keras, and TensorFlow} (3rd ed.). O'Reilly Media. — Abordagem prática que complementa este livro.
    \item Bradski G, Kaehler A. (2008). \textit{Learning OpenCV}. O'Reilly Media. — Referência clássica para processamento de imagens.
    \item Mason R (2022). \textit{DICOM is Easy} (3rd ed.). — Guia prático para trabalhar com imagens DICOM em Python usando pydicom.
\end{enumerate}

\subsection{Pontos-Chave}
\begin{itemize}
    \item Python é a linguagem franca da IA odontológica devido ao seu ecossistema científico maduro, curva de aprendizado suave e acesso gratuito a GPU via Google Colab.
    \item NumPy arrays são a representação computacional natural de imagens radiográficas (2D), volumes CBCT (3D) e matrizes de características clínicas.
    \item Pandas DataFrames estruturam dados clínicos odontológicos (periodontais, CPOD, prontuários) permitindo análises estatísticas com poucas linhas de código.
    \item Matplotlib + Seaborn transformam dados brutos em visualizações publicáveis que revelam padrões clínicos.
    \item O protocolo SDD+TDD garante que notebooks odontológicos sejam não apenas corretos, mas também verificáveis e reprodutíveis.
\end{itemize}

\subsection{Questões de Revisão}
\begin{enumerate}
    \item Como você representaria um exame periodontal completo (168 sítios, 6 faces por dente) usando Pandas?
    \item Qual a vantagem de usar CLAHE em vez de equalização de histograma global em radiografias odontológicas?
    \item Por que a documentação de notebooks com SDD+TDD é importante para a reprodutibilidade científica?
    \item Como o NumPy permite realizar operações em imagens radiográficas sem loops explícitos em Python?
    \item Que tipo de GPU o Google Colab oferece gratuitamente e qual sua relevância para treinamento de CNNs odontológicas?
\end{enumerate}

%==============================================================================
% FIM DO CAPÍTULO 3
%==============================================================================
